\section{Attack Surface 1: Census P.L.\ 94-171 Population Counts}
\label{sec:pl94}

\subsection{The Data Format}

The Census P.L.\ 94-171 redistricting data file provides population counts
for every census geographic unit from the state level to the block level.
For \bisect{}, the relevant unit is the census tract (GEOID prefix ending
in ``00'', indicating a 2020 census tract).
Each tract record contains:
\begin{itemize}
\item Total population (P1\_001N)
\item Population by race (P1\_002N through P1\_009N)
\item Hispanic or Latino origin (P2\_002N)
\item Voting-age population by race (P3, P4 tables)
\end{itemize}

The algorithm uses P1\_001N (total population) for the equal-population
constraint and the P3/P4 tables for VRA boost weight computation.

\subsection{The Attack}

A hostile actor who controls the P.L.\ 94-171 file could systematically:
\begin{enumerate}
\item Inflate population counts in tracts that vote heavily Republican
  (rural areas, outer suburbs).
\item Deflate population counts in tracts that vote heavily Democratic
  (urban cores, college towns).
\end{enumerate}

This distortion biases the equal-population constraint: the algorithm
will split tracts into districts that are ``equal'' in the manipulated
population counts but systematically different in true population.
The practical effect is that Republican-heavy tracts become overrepresented
in the sense that the algorithm assigns fewer such tracts per district
than warranted by their true population.

\subsection{Effect Magnitude}

The key constraint on this attack is averaging.
A $\pm 1\%$ manipulation per tract produces a distortion of $\pm 0.01 \times
p_i$ people in tract $i$ (where $p_i$ is the true population).
For an average-sized congressional district ($\sim$760,000 people, drawn
from $\sim$18 tracts in a 14-district state):
\[
  \text{District-level distortion} = \pm 1\% \times \frac{p_{\rm state}}{k}
  \approx \pm 7{,}600 \text{ people}.
\]

To flip one competitive district from Democratic to Republican using
population manipulation alone, the adversary must bias the boundaries
enough to systematically include Republican-heavy tracts at the expense
of Democratic-heavy tracts.
Given that competitive districts are typically within 2--3\% of the
partisan threshold ($D$-vote share $\in [47\%, 53\%]$), and given that
the algorithm adds $\pm 7{,}600$ people per district through manipulation,
the expected vote-share change in a competitive district is approximately:
\[
  \Delta D\text{-share per district} \approx \pm 1\% \times R = \pm 0.3\%,
\]
where $R$ is the ratio of partisan correlation with the manipulation to
the district size.
This is far below the 3--6\% margin of a competitive district.

Nationally, summing across all competitive districts and applying the
saturation argument from R.0~\citep{r0paper}: $\Delta D \leq 0.4$ seats.

\subsection{Detection}

\texttt{bisect fetch} computes:
\begin{verbatim}
SHA256(pl94171_state_XX_2020.pl) -> stored in manifest
\end{verbatim}
and compares against the Census Bureau's published checksum at:
\begin{verbatim}
https://www2.census.gov/programs-surveys/decennial/2020/
  data/01-Redistricting_File--PL_94-171/
  [StateXX]/2020.pl.sha256
\end{verbatim}

Any byte-level modification to the P.L.\ 94-171 file changes the SHA-256
hash with probability $1 - 2^{-256}$ (i.e., with certainty for any
practical computation).
